\documentclass[12pt]{article}
\usepackage[utf8]{inputenc}
\usepackage[T1]{fontenc}
\usepackage[margin=1in]{geometry}
\usepackage{setspace}
\usepackage{xcolor}
\usepackage[colorlinks=true, urlcolor=blue!60!black, linkcolor=blue!60!black]{hyperref}
\setlength{\parindent}{2em}
\setlength{\parskip}{1.25em}
\renewcommand{\baselinestretch}{1.0}
\title{\textit{U.S. Department of Agriculture v. Moreno}, 413 U.S. 528 (1973)}
\author{Case Brief}
\date{}
\begin{document}
\maketitle

\section{Case Metadata}

\begin{itemize}
  \item \textbf{Case Name:} \href{https://supreme.justia.com/cases/federal/us/413/528/}{\textit{United States Department of Agriculture v. Moreno}}
  \item \textbf{Citation:} 413 U.S. 528 (1973); Docket No. 72-534 (\href{https://www.law.cornell.edu/supremecourt/text/413/528}{Cornell LII}; \href{https://tile.loc.gov/storage-services/service/ll/usrep/usrep413/usrep413528/usrep413528.pdf}{Library of Congress U.S. Reports PDF})
  \item \textbf{Dates:} Argued April 23, 1973; Decided June 25, 1973
  \item \textbf{Court Below:} Three-Judge United States District Court for the District of Columbia, 345 F. Supp. 310 (D.D.C. 1972) (direct appeal)
  \item \textbf{Disposition:} Affirmed (7--2)
  \item \textbf{Vote and Alignment:}
  \begin{itemize}
    \item \textbf{Majority (7):} Justice William J. Brennan, Jr. (author), joined by Justices William O. Douglas, Potter Stewart, Byron R. White, Thurgood Marshall, Harry A. Blackmun, and Lewis F. Powell, Jr.
    \item \textbf{Concurrence (1):} Justice William O. Douglas (concurring opinion, 413 U.S. at 538--545).
    \item \textbf{Dissent (2):} Justice William H. Rehnquist (author), joined by Chief Justice Warren E. Burger (413 U.S. at 545--547).
  \end{itemize}
  \item \textbf{Constitutional Provision \& Statute:} Fifth Amendment Due Process Clause (equal protection component via \textit{Bolling v. Sharpe}, 347 U.S. 497 (1954)); Food Stamp Act of 1964, \S~3(e), 7 U.S.C. \S~2012(e), as amended in 1971 (84 Stat. 2048); 7 C.F.R. \S~270.2(jj).
\end{itemize}

\section{Facts and Procedural History}

Congress created the federal food stamp program in 1964 to ``alleviate hunger and malnutrition'' among low-income households and strengthen the agricultural economy (7 U.S.C. \S~2011). Originally, \S~3(e) defined an eligible ``household'' as ``a group of related or non-related individuals'' living as one economic unit sharing cooking facilities and purchasing food in common (\href{https://supreme.justia.com/cases/federal/us/413/528/}{413 U.S. at 530}). In January 1971, however, a conference committee added---without committee hearings---an amendment redefining ``household'' to exclude any group containing an individual unrelated to any other member (with narrow exceptions for persons over age 60 and certain live-in attendants). Scant legislative history showed the amendment was aimed at excluding ``hippies'' and ``hippie communes'' from food stamps (H.R. Conf. Rep. No. 91-1793, at 8; 116 Cong. Rec. 44439 (1970)) (\href{https://supreme.justia.com/cases/federal/us/413/528/}{413 U.S. at 530--531, 534 \& n.6}).

In practice, the exclusion fell harshly on destitute households doubling up to survive. Appellee Jacinta Moreno, a 56-year-old diabetic receiving \$75 per month in public assistance, lived with Ermina Sanchez and Sanchez's three children to share expenses and receive care; after rent, utilities, hospital transit, and laundry, Moreno had \$10 per month left for food, yet both she and the Sanchez family faced total disqualification solely because they were unrelated. Appellee Sheilah Hejny and her indigent family faced cutoff for sheltering an emotionally troubled 20-year-old woman. Appellee Victoria Keppler lost food stamps because she shared an apartment near a school for the deaf with an unrelated woman on public assistance so her hearing-impaired daughter could attend (\href{https://supreme.justia.com/cases/federal/us/413/528/}{413 U.S. at 531--532}).

Moreno and other plaintiffs filed a class action against the U.S. Department of Agriculture (USDA) in the District of Columbia. A three-judge District Court held the ``unrelated person'' exclusion unconstitutional under the Fifth Amendment's equal protection component (345 F. Supp. 310 (D.D.C. 1972)), and the USDA appealed directly to the Supreme Court. Decided the same day as its companion due-process case, \href{https://supreme.justia.com/cases/federal/us/413/508/}{\textit{U.S. Department of Agriculture v. Murry}, 413 U.S. 508 (1973)}, \textit{Moreno} tested the outer limits of rational-basis scrutiny.

\section{Issue(s) Presented}

Whether \S~3(e) of the Food Stamp Act of 1964, as amended in 1971 to exclude from federal food assistance any otherwise eligible household containing an individual unrelated to any other household member, violates the equal protection component of the Fifth Amendment's Due Process Clause under rational-basis review.

\section{Holding}

\textbf{Yes; judgment affirmed (7--2).} Writing for the Court, Justice Brennan held that \S~3(e)'s ``unrelated person'' classification violated Fifth Amendment equal protection even under traditional rational-basis scrutiny (\href{https://supreme.justia.com/cases/federal/us/413/528/}{413 U.S. at 533, 538}). The classification was irrelevant to the Act's nutritional and agricultural purposes; its actual legislative purpose---barring ``hippies'' and ``hippie communes''---was constitutionally illegitimate because ``a bare congressional desire to harm a politically unpopular group cannot constitute a legitimate governmental interest''; and the Government's post hoc fraud-prevention rationale failed because the provision was ``wholly without any rational basis'' (\href{https://supreme.justia.com/cases/federal/us/413/528/}{413 U.S. at 534--538}).

\section{Reasoning}

\subsection{Irrelevance to Statutory Purposes and Rejection of Legislative Animus}

Applying Fifth Amendment reverse-incorporation equal protection (\textit{Bolling v. Sharpe}; \textit{Schneider v. Rusk}, 377 U.S. 163, 168 (1964)), Justice Brennan nominally invoked ordinary rational-basis review: a classification must be ``rationally related to a legitimate governmental interest'' (\href{https://supreme.justia.com/cases/federal/us/413/528/}{413 U.S. at 533}). First, familial relationships had nothing to do with the Act's express goals in 7 U.S.C. \S~2011---raising low-income nutrition and stimulating agricultural demand. Second, the sparse legislative history revealed an actual purpose to exclude ``hippies'' and ``hippie communes.'' Brennan declared that purpose flatly illegitimate: equal protection forbids legislation grounded in a ``bare congressional desire to harm a politically unpopular group'' (\href{https://supreme.justia.com/cases/federal/us/413/528/}{413 U.S. at 534--535}). (The Government abandoned its District Court ``morality'' defense on appeal.)

\subsection{Skeptical Scrutiny of the Fraud-Prevention Rationale (``Rational Basis with Bite'')}

The Government pivoted to a facially legitimate objective: minimizing program fraud, arguing that unrelated households are more likely to include voluntarily poor abusers and are ``relatively unstable,'' making fraud harder to detect. Departing sharply from the extreme deference of \textit{Dandridge v. Williams}, 397 U.S. 471 (1970), Brennan dismantled this rationale on two empirical and structural grounds (\href{https://supreme.justia.com/cases/federal/us/413/528/}{413 U.S. at 535--538}):
\begin{enumerate}
  \item \textbf{Existing Targeted Anti-Fraud Provisions:} The Food Stamp Act already contained dedicated provisions disqualifying able-bodied adults who refuse work (\S~5(c), 7 U.S.C. \S~2014(c)) and imposing felony criminal penalties for food stamp fraud (\S\S~14(b)--(c), 7 U.S.C. \S\S~2023(b)--(c)). Those targeted tools cast ``considerable doubt'' on the claim that \S~3(e) was rationally aimed at the same abuses.
  \item \textbf{Perverse Practical Operation:} Because \S~3(e) defined a ``household'' as persons living as one economic unit, sharing cooking facilities, and buying food in common, any savvy ``hippie'' or would-be abuser could easily evade the ban by splitting food purchases into two separate ``households.'' Citing California's Director of Social Welfare, Brennan concluded that in practical operation \S~3(e) excluded \textit{only} the most desperately poor households---such as Moreno, Hejny, and Keppler---who could not afford to rearrange their shared living arrangements. The classification was therefore not merely imprecise, but ``wholly without any rational basis'' (\href{https://supreme.justia.com/cases/federal/us/413/528/}{413 U.S. at 537--538}).
\end{enumerate}

\section{Concurring and Dissenting Opinions}

\begin{itemize}
  \item \textbf{Justice Douglas (Concurring, 413 U.S. at 538--545):} Douglas candidly conceded that, absent a fundamental right, he ``could not say that this `unrelated' person provision has no `rational' relation to control of fraud,'' and \textit{Dandridge} would otherwise control. However, Douglas argued that \S~3(e) penalized the First Amendment \textbf{freedom of association}---``banding together to combat the common foe of hunger''---triggering strict scrutiny (requiring a ``compelling'' interest and narrow tailoring under \textit{NAACP v. Alabama} and \textit{Shapiro v. Thompson}), which \S~3(e) failed (\href{https://supreme.justia.com/cases/federal/us/413/528/}{413 U.S. at 541--545}).
  \item \textbf{Justice Rehnquist (Dissenting, joined by Chief Justice Burger, 413 U.S. at 545--547):} Rehnquist accused the majority of writing ``a very persuasive congressional committee report'' rather than applying deferential rational-basis review. Under \textit{Dandridge} and \textit{Mourning v. Family Publications Service, Inc.}, 411 U.S. 356 (1973), Congress could reasonably decide to subsidize only related family units because blood or marital ties provide a rough assurance that the household was not formed solely to collect food stamps. That Congress used a ``blunt instrument'' with ``unfortunate and perhaps unintended consequences'' was a policy flaw for Congress, not a constitutional defect (\href{https://supreme.justia.com/cases/federal/us/413/528/}{413 U.S. at 545--547}).
\end{itemize}

\section{Key Quotations}

\begin{enumerate}
  \item \textbf{The Anti-Animus Principle (Brennan, J., \href{https://supreme.justia.com/cases/federal/us/413/528/}{413 U.S. at 534}):}
  \begin{quote}
    ``For if the constitutional conception of `equal protection of the laws' means anything, it must, at the very least, mean that a bare congressional desire to harm a politically unpopular group cannot constitute a legitimate governmental interest.''
  \end{quote}

  \item \textbf{Practical Operation and Irrationality (Brennan, J., \href{https://supreme.justia.com/cases/federal/us/413/528/}{413 U.S. at 538}):}
  \begin{quote}
    ``Thus, in practical operation, the 1971 amendment excludes from participation in the food stamp program, not those persons who are `likely to abuse the program' but, rather, only those persons who are so desperately in need of aid that they cannot even afford to alter their living arrangements so as to retain their eligibility.''
  \end{quote}

  \item \textbf{Freedom of Association in the Home (Douglas, J., concurring, \href{https://supreme.justia.com/cases/federal/us/413/528/}{413 U.S. at 543}):}
  \begin{quote}
    ``The right of association, the right to invite the stranger into one's home is too basic in our constitutional regime to deal with roughshod.''
  \end{quote}

  \item \textbf{Deferential Rational-Basis Critique (Rehnquist, J., dissenting, \href{https://supreme.justia.com/cases/federal/us/413/528/}{413 U.S. at 545--546}):}
  \begin{quote}
    ``The Court's opinion would make a very persuasive congressional committee report arguing against the adoption of the limitation in question. Undoubtedly, Congress attacked the problem with a rather blunt instrument . . . . But questions such as this are for Congress, rather than for this Court . . . .''
  \end{quote}
\end{enumerate}

\section{Doctrinal Significance and Modern Legacy (Through 2026)}

\textit{Moreno} is the foundational cornerstone of two interlocking equal-protection doctrines: \textbf{``rational basis with bite''} and the \textbf{constitutional anti-animus principle}. Together with three later landmarks, it forms the Court's ``animus quadrilogy'':
\begin{itemize}
  \item \textbf{The Animus Quadrilogy:} The Court directly invoked \textit{Moreno}'s anti-animus rule to invalidate a zoning denial targeting a group home for intellectually disabled residents in \href{https://supreme.justia.com/cases/federal/us/473/432/}{\textit{City of Cleburne v. Cleburne Living Center, Inc.}, 473 U.S. 432, 446--447 (1985)}; Colorado's Amendment 2 barring antidiscrimination protections for gay and lesbian citizens in \href{https://supreme.justia.com/cases/federal/us/517/620/}{\textit{Romer v. Evans}, 517 U.S. 620, 634--635 (1996)}; and Section 3 of the Defense of Marriage Act (DOMA) under Fifth Amendment equal protection in \href{https://supreme.justia.com/cases/federal/us/570/744/}{\textit{United States v. Windsor}, 570 U.S. 744, 770 (2013)}. Justice O'Connor also synthesized \textit{Moreno}, \textit{Cleburne}, and \textit{Romer} in her equal-protection concurrence in \href{https://supreme.justia.com/cases/federal/us/539/558/}{\textit{Lawrence v. Texas}, 539 U.S. 558, 580 (2003)}, observing that the Court applies a ``more searching form of rational basis review'' whenever a law exhibits a desire to harm an unpopular group.
  \item \textbf{Retrenchment and Modern Status Through 2026:} In the Roberts Court era, the animus doctrine has narrowed considerably. In \href{https://supreme.justia.com/cases/federal/us/585/17-965/}{\textit{Trump v. Hawaii}, 585 U.S. 667 (2018)}, the Court upheld the President's travel proclamation despite challenged statements, stressing that a policy survives rational-basis review unless it is ``inexplicable by anything but animus.'' Most recently, in \href{https://www.supremecourt.gov/opinions/24pdf/23-477_2cp3.pdf}{\textit{United States v. Skrmetti}, 605 U.S. 495 (2025)}, the Court upheld Tennessee's ban on gender-affirming medical care for minors under deferential rational-basis review (\textit{FCC v. Beach Communications, Inc.}, 508 U.S. 307 (1993)) without finding animus, even as lower federal courts through 2025--2026 continue to invoke \textit{Moreno}'s animus doctrine in challenges to federal executive policies affecting transgender individuals (see \href{https://scholarship.law.tamu.edu/cgi/viewcontent.cgi?article=1438&context=lawreview}{Texas A\&M L. Rev. (Oct. 2026)}). Thus, through 2026, \textit{Moreno} remains good law---and the primary textual hook for challenging classifications targeting non-suspect groups---yet courts rarely invalidate laws under \textit{Moreno} if the government offers any plausible, independent regulatory justification.
\end{itemize}

\section{Pedagogical Critiques and Exam Hypotheticals}

\subsection{Key Pedagogical Critiques}
\begin{itemize}
  \item \textbf{Doctrinal Candor vs. Tiered Scrutiny:} Critics (including Justice Rehnquist in dissent and Justice Douglas in concurrence) note that \textit{Moreno} purports to apply ordinary rational-basis review while covertly deploying heightened means-ends scrutiny. Because the Court refused to declare poverty, non-traditional households, or hippies a suspect class, it left lower courts without clear rules for when to apply \textit{Beach Communications} deference versus \textit{Moreno} ``bite.''
  \item \textbf{Proving Legislative Animus:} Relying on stray floor remarks (such as Senator Holland's anti-hippie comment) raises classic motive-inquiry problems: how many legislators must harbor animus, and does any independent, non-pretextual rationale (like fraud prevention) cleanse a statute enacted partly out of hostility?
\end{itemize}

\subsection{Exam Hypotheticals (Including ``Opposite Side'' Contrast)}

\subsubsection{Hypothetical 1 (Classic \textit{Moreno} / \textit{Cleburne} Animus Application)}
Facing complaints from homeowners about noise and parking near a university, a city amends its municipal housing-voucher program to disqualify any household containing ``an individual with a prior misdemeanor conviction for participating in an unpermitted street protest.'' During floor debate, the ordinance's sponsor declares, ``Taxpayers shouldn't subsidize radical campus agitators.'' When sued, the city argues the rule rationally conserves limited housing funds and promotes neighborhood stability.

\textbf{Analysis:} Because protest misdemeanants are not a suspect class, rational-basis review applies (alongside First Amendment retaliation concerns). Under \href{https://supreme.justia.com/cases/federal/us/413/528/}{\textit{Moreno}}, the sponsor's statement flags a ``bare desire to harm a politically unpopular group.'' Like the fraud rationale in \textit{Moreno}, the city's post hoc justifications fail ``rational basis with bite'': violent felons and noisy tenants remain eligible while peaceful misdemeanor protesters are singled out, rendering the means-ends fit grossly under- and over-inclusive and inexplicable by anything other than animus.

\subsubsection{Hypothetical 2 (``Opposite Side'' Contrast---Invalid Animus vs. Valid Line-Drawing)}
\begin{itemize}
  \item \textbf{Law A (Invalid Under \textit{Moreno}):} A state excludes from emergency utility-bill relief any otherwise eligible household that includes an unrelated adult who works as a tattoo artist or body piercer, following legislative remarks denouncing ``counterculture freeloaders.'' Even if the state argues cash-based trades underreport income, existing income-verification rules address fraud, and targeting one stigmatized occupation fails \textit{Moreno}.
  \item \textbf{Law B (``Opposite Side''---Valid Under Deferential Rational Basis):} A town enacts a neutral residential zoning ordinance limiting single-family zones to related persons or no more than two unrelated persons living together, citing traffic congestion, noise, and neighborhood quiet (modeled on \href{https://supreme.justia.com/cases/federal/us/416/1/}{\textit{Village of Belle Terre v. Boraas}, 416 U.S. 1 (1974)}). Unlike \textit{Moreno}, Law B lacks record evidence of bare animus toward an unpopular group, does not strip subsistence aid from destitute households, and bears a plausible, non-pretextual connection to legitimate land-use goals---surviving traditional rational-basis review.
\end{itemize}

\end{document}
